\chapter{Introduction}
\section{Motivation}

University websites contain a large amount of important but fragmented information, including academic announcements, administrative notices, and campus services. However, traditional keyword-based search systems are often insufficient for accurately understanding user intent, which makes information retrieval inefficient and inconvenient.

With the development of large language models (LLMs), Retrieval-Augmented Generation (RAG) has emerged as an effective approach to combine external knowledge retrieval with generative capabilities. This enables systems to provide more accurate, context-aware, and natural language responses.

Motivated by these limitations and opportunities, this project aims to build a local-first RAG system tailored for SUSTech's public web content, improving the accessibility and usability of campus information through natural language interaction.


 \section{Contributions}\label{sec:motiv}
 This project makes the following contributions:

\begin{itemize}
    \item We design and implement a complete Retrieval-Augmented Generation (RAG) pipeline for SUSTech public web content, integrating data acquisition, preprocessing, retrieval, reranking, and generation.

    \item We develop a modular backend system based on FastAPI, supporting multiple LLM inference backends including llama.cpp and vLLM, enabling flexible deployment across different hardware environments.

    \item We construct a semantic retrieval system using BGE embedding and reranking models, improving the relevance of retrieved context for generation.

    \item We design a Vue 3-based frontend interface with streaming response support, multi-session management, and customizable UI, providing an interactive user experience.

    \item We build a local-first architecture that allows the system to operate independently of external APIs, ensuring privacy, controllability, and extensibility.
\end{itemize}

\section{System Overview}

The proposed system follows a standard Retrieval-Augmented Generation (RAG) pipeline, consisting of both backend and frontend components.

The backend is responsible for data processing and model inference. It first collects publicly available web data from the SUSTech official website. The raw text is then cleaned, segmented into chunks, and transformed into vector embeddings using BGE models. These embeddings are stored in a vector index for efficient similarity search. During query time, the system retrieves relevant chunks, applies a reranking model to improve relevance, and finally feeds the optimized context into a large language model (Qwen3) for response generation.

The frontend provides an interactive user interface built with Vue 3 and Vite. It communicates with the backend through RESTful APIs and supports real-time streaming responses using Server-Sent Events (SSE). Users can interact with the system through a web-based chat interface, manage multiple sessions, and customize the UI appearance.

Overall, the system forms a complete end-to-end RAG workflow that connects web data acquisition, semantic retrieval, and generative AI into a unified application.


\chapter{System Architecture}


\section{Architecture Overview}

The proposed system follows a modular Retrieval-Augmented Generation (RAG) architecture, consisting of three main components: data ingestion, backend RAG pipeline, and frontend interaction interface. The overall goal is to transform unstructured web content from the Southern University of Science and Technology (SUSTech) official website into a queryable semantic knowledge system.

The system is designed with a clear separation between data processing, retrieval and generation, and user interaction layers. This modular design improves maintainability, scalability, and deployment flexibility.

\section{Backend and Frontend Interaction}

The system is divided into a backend service and a frontend application.

The backend is implemented using Python and FastAPI. It is responsible for document processing, embedding generation, vector indexing, retrieval, reranking, and large language model inference. The backend exposes RESTful APIs that handle query requests and return streaming responses.

The frontend is built using Vue 3 and Vite. It provides an interactive chat-based interface that communicates with the backend through HTTP requests and receives responses via Server-Sent Events (SSE). The frontend also supports multi-session management, local caching of conversations, and UI customization.

This separation of concerns allows the system to independently scale backend computation and frontend interaction layers.

\section{Design Characteristics}

The proposed architecture has the following key characteristics:

\begin{itemize}
    \item \textbf{Modularity:} Each component of the system (ingestion, retrieval, generation, UI) is independently designed and can be replaced or upgraded without affecting the overall system.

    \item \textbf{Flexibility:} The backend supports multiple LLM inference engines, including llama.cpp and vLLM, enabling deployment across different hardware environments.

    \item \textbf{Scalability:} Vector-based retrieval allows efficient search over large-scale document collections.

    \item \textbf{Interactivity:} Streaming responses via SSE enable real-time user interaction.
\end{itemize}
\chapter{Backend Design}
\section{RAG Pipeline Overview}

The core workflow of the system follows a standard Retrieval-Augmented Generation (RAG) pipeline designed for processing publicly available web content from the SUSTech official website.

The pipeline consists of the following sequential stages:

Public web pages are first collected through a crawling module, which extracts raw HTML content and stores it in a structured format (raw\_documents.jsonl). The raw data is then passed through a preprocessing stage, where noise such as navigation elements, scripts, and redundant markup is removed, producing cleaned documents (documents.cleaned.jsonl).

Next, the cleaned documents are segmented into smaller semantic units through a chunking process, generating chunks.jsonl. These chunks are then embedded and indexed into a vector database (ChromaDB), enabling efficient similarity-based retrieval.

During inference, user queries are first encoded and used to retrieve the most relevant chunks. A reranking model is then applied to refine retrieval results by improving semantic relevance.

Finally, the optimized context is passed into a large language model (llama.cpp or vLLM backend), which generates a response in a streaming manner to the frontend interface.

\section{Pipeline Description}

The complete system workflow is illustrated as follows:

\begin{itemize}
    \item Public Web Pages
    \item Crawling $\rightarrow$ raw\_documents.jsonl + HTML
    \item Preprocessing $\rightarrow$ documents.cleaned.jsonl
    \item Chunking $\rightarrow$ chunks.jsonl
    \item Indexing $\rightarrow$ ChromaDB
    \item Retrieval + Reranking
    \item LLM Generation (llama.cpp / vLLM)
    \item Streaming Output to Frontend
\end{itemize}

\chapter{Frontend Routing & Integration}
\section{Routing and Page Structure}

The frontend application defines multiple entry routes to support different usage scenarios, including desktop, mobile, embedded, and auxiliary tool modes. Each route corresponds to a specific UI layout optimized for its intended environment.

The routing structure is as follows:

\begin{itemize}
    \item \textbf{/} — Desktop chat interface
    \item \textbf{/mobile} — Mobile-optimized chat interface
    \item \textbf{/settings} — Configuration and settings page
    \item \textbf{/embed} — Full-page embedded interface
    \item \textbf{/ball} — Floating entry widget
    \item \textbf{/binarize.html} — Standalone static utility page
\end{itemize}

The main application component (App.vue) dynamically switches between desktop and mobile layouts based on the browser window width. However, embedded, floating, and settings pages are excluded from this automatic layout switching mechanism to preserve their specific UI constraints.

\section{Backend Integration}

For system integration, the backend service is first started using the following command:

\begin{verbatim}
uv run sustech-rag serve --host 127.0.0.1 --port 8001
\end{verbatim}

The frontend is then launched separately using:

\begin{verbatim}
npm run dev
\end{verbatim}

By default, the frontend communicates with the backend through the relative API base path \texttt{/api}. In cases where the frontend and backend are deployed on different origins, the API base URL can be explicitly configured to a full endpoint such as:

\begin{verbatim}
http://127.0.0.1:8001/api
\end{verbatim}

The backend CORS policy is configured to allow requests from the following origins:

\begin{itemize}
    \item http://127.0.0.1:3000
    \item http://localhost:3000
\end{itemize}

\section{State Management and Persistence}

The frontend uses browser \texttt{localStorage} for persistent state management. Two primary data structures are stored locally:

\begin{itemize}
    \item Chat history: \texttt{ragwebui:chats:v1}
    \item User settings: \texttt{ragwebui:settings:v1}
\end{itemize}

Upon initial load, the frontend sends a request to \texttt{POST /api/identity} to obtain a user identity ID. If the backend service is unavailable, a local fallback ID is generated to ensure continued operation in offline or degraded modes.

\section{SSE Event Protocol}

The frontend supports a structured Server-Sent Events (SSE) protocol for streaming model responses. The supported event types include:

\begin{itemize}
    \item start
    \item reference
    \item think.delta
    \item think.end
    \item content.delta
    \item tool.call
    \item tool.result
    \item image
    \item error
    \item done
\end{itemize}

In the current backend implementation, the primary emitted events include \texttt{start}, \texttt{reference}, \texttt{think.delta}, \texttt{think.end}, \texttt{content.delta}, \texttt{error}, and \texttt{done}.

A detailed specification of the communication protocol is documented in \texttt{docs/API.md}, with an additional browser-accessible mirrored version available at \texttt{public/API.md}.


\chapter{Conclusion}

This report presents the design and implementation of a local-first Retrieval-Augmented Generation (RAG) system for querying publicly available web content from the Southern University of Science and Technology (SUSTech).

The system integrates a complete pipeline including web data acquisition, preprocessing, semantic chunking, vector-based retrieval, reranking, and large language model generation. The backend is built with a modular architecture supporting multiple inference engines such as llama.cpp and vLLM, while the frontend provides a responsive and interactive interface based on Vue 3 and Vite with real-time streaming support.

By combining retrieval-based methods with generative models, the system significantly improves the efficiency and usability of accessing campus information. The modular design ensures flexibility, scalability, and ease of deployment across different environments.

Future work may include improving retrieval accuracy, expanding data sources, and integrating more advanced reasoning capabilities in the generation model.


